\documentclass[12pt, a4paper]{article}

\usepackage[utf8]{inputenc}
\usepackage[T1]{fontenc}
\usepackage[spanish]{babel}
\usepackage{geometry}
\geometry{margin=2.5cm}
\usepackage{booktabs}
\usepackage{tabularx}
\usepackage{array}
\usepackage{listings}
\usepackage{xcolor}
\usepackage{titlesec}
\usepackage{parskip}
\usepackage{hyperref}
\usepackage{amsmath}

% Colores
\definecolor{codebg}{RGB}{245,245,245}
\definecolor{codeframe}{RGB}{200,200,200}
\definecolor{keywordcolor}{RGB}{0,0,180}
\definecolor{commentcolor}{RGB}{80,130,80}
\definecolor{stringcolor}{RGB}{160,0,0}
\definecolor{accentblue}{RGB}{30,90,160}
\definecolor{accentgreen}{RGB}{30,130,80}
\definecolor{accentred}{RGB}{160,40,40}
\definecolor{tabheader}{RGB}{230,237,245}

% Estilo de código Java
\lstdefinestyle{java}{
  language=Java,
  backgroundcolor=\color{codebg},
  frame=single,
  rulecolor=\color{codeframe},
  basicstyle=\ttfamily\small,
  keywordstyle=\color{keywordcolor}\bfseries,
  commentstyle=\color{commentcolor}\itshape,
  stringstyle=\color{stringcolor},
  breaklines=true,
  breakatwhitespace=true,
  showstringspaces=false,
  tabsize=4,
  numbers=none,
  xleftmargin=8pt,
  xrightmargin=8pt,
  aboveskip=6pt,
  belowskip=6pt,
}

% Títulos
\titleformat{\section}{\large\bfseries\color{accentblue}}{}{0em}{}[\titlerule]
\titleformat{\subsection}{\normalsize\bfseries\color{accentblue}}{}{0em}{}

\hypersetup{
  colorlinks=true,
  linkcolor=accentblue,
  urlcolor=accentblue
}

\title{
  \textbf{Métodos de la clase Vertex}\\[0.4em]
  \large Algoritmos A* y Backtracking --- ADDA v6\\[0.2em]
  \normalsize PI4\_2526\_PBase
}
\author{}
\date{\today}

\begin{document}

\maketitle
\tableofcontents
\newpage

% ============================================================
\section{Contexto general}
% ============================================================

En el framework ADDA v6, cada ejercicio modela el espacio de búsqueda como un \textbf{grafo virtual}. Los nodos de ese grafo son instancias de una clase \texttt{Vertex} que implementa la interfaz \texttt{VirtualVertex<V, E, Action>}.

Tanto A* como Backtracking (BT) recorren ese grafo. La clase \texttt{Vertex} se implementa \textbf{una sola vez} y sirve para ambos algoritmos; la diferencia está en qué métodos extra se añaden fuera del vértice:

\begin{itemize}
  \item \textbf{A*} requiere una clase de \textbf{heurística externa} que estime el coste restante.
  \item \textbf{BT con semilla greedy} requiere el método \texttt{greedyAction()} dentro del vértice.
\end{itemize}

El flujo general de decisión es:

\[
\texttt{start()} \;\longrightarrow\; \texttt{actions()} \;\longrightarrow\; \texttt{neighbor(a)} \;\longrightarrow\; \texttt{isValid()?} \;\longrightarrow\; \texttt{goal()?} \;\longrightarrow\; \texttt{goalHasSolution()}
\]

% ============================================================
\section{Métodos para A*}
% ============================================================

A* construye el grafo virtual y expande nodos ordenados por $f(n) = g(n) + h(n)$, donde $g(n)$ es el coste acumulado y $h(n)$ es la estimación heurística del coste restante.

\subsection{actions()}

Devuelve la lista de \textbf{decisiones posibles} desde el estado actual. A* las usa para generar los nodos vecinos a explorar. Si no hay acciones (o ya se alcanzó el goal), devuelve lista vacía.

\begin{lstlisting}[style=java]
// Vertex2I -- en que contenedor puedo meter el elemento actual?
public List<Integer> actions() {
    if (this.indice().equals(Datos2.getNumElementos())) return List.of();
    return IntStream.range(0, Datos2.getNumContenedores()).boxed()
        .filter(x -> !contenedoresCompletos().contains(x))
        .filter(x -> Datos2.getPuedeUbicarse(this.indice(), x))
        .filter(x -> Datos2.getTamElemento(this.indice()) <= capRest().get(x))
        .toList();
}
\end{lstlisting}

\subsection{neighbor(action)}

Dado que se toma la acción \texttt{a}, devuelve el \textbf{nuevo estado} resultante. Aquí se actualiza todo lo que cambia al tomar esa decisión. El campo \texttt{indice + 1} indica que pasamos a la siguiente decisión.

\begin{lstlisting}[style=java]
// Vertex2I -- asigna el elemento actual al contenedor a
public Vertex2I neighbor(Integer a) {
    List<Integer> capRest2 = List2.copy(tamanoContenedores);
    capRest2.set(a, capRest2.get(a) - Datos2.getTamElemento(indice));
    return Vertex2I.of(indice + 1, capRest2, ...);
}
\end{lstlisting}

\subsection{edge(action)}

Crea el objeto \textbf{arista} entre el vértice actual y su vecino. Es un wrapper sobre \texttt{neighbor}. El algoritmo lo usa para construir el \texttt{GraphPath} y calcular $g(n)$.

\begin{lstlisting}[style=java]
// Vertex2I
public Edge2 edge(Integer a) {
    return Edge2.of(this, this.neighbor(a), a);
}
\end{lstlisting}

\subsection{goal()}

Devuelve \texttt{true} cuando \textbf{ya no quedan decisiones} por tomar. No indica si la solución es válida, solo que el árbol de decisión ha terminado.

\begin{lstlisting}[style=java]
// Vertex2I -- se han procesado todos los elementos
public Boolean goal() {
    return this.indice() == Datos2.getNumElementos();
}

// Vertex3I -- se han visitado todos los nodos del grafo
public Boolean goal() {
    return this.indice == Datos3.N;
}
\end{lstlisting}

\subsection{goalHasSolution()}

Solo se invoca cuando \texttt{goal()} es \texttt{true}. Comprueba si el estado final \textbf{cumple todas las restricciones}, es decir, si lo construido es una solución válida.

\begin{lstlisting}[style=java]
// Vertex3I -- llego al final, visito >=2 monumentos consecutivos y no supera el tiempo
public Boolean goalHasSolution() {
    return this.indice == Datos3.N
        && this.nConsec >= 2
        && this.durAcum >= 0.0;
}
\end{lstlisting}

\subsection{isValid()}

Poda anticipada: se evalúa \textbf{antes de explorar} un nodo. Si devuelve \texttt{false}, ese nodo y toda su subrama se descartan. Cuanto más restrictivo y barato sea este método, más eficiente es el algoritmo.

\begin{lstlisting}[style=java]
// CandidatosHyperVertex -- presupuesto no negativo e indice en rango
public Boolean isValid() {
    return this.indice >= 0
        && this.indice <= Datos1.getNumCandidatos()
        && this.presRest <= Datos1.getPresupuestoMax();
}
\end{lstlisting}

\subsection{Heurística (clase externa) --- exclusivo A*}

Estima de forma \textbf{optimista} cuánto queda por ganar/perder hasta el final. Debe ser \textbf{admisible}: nunca sobreestimar si se minimiza, nunca subestimar si se maximiza. A* la usa para ordenar qué nodo explorar primero mediante $f(n) = g(n) + h(n)$.

\begin{lstlisting}[style=java]
// Heuristic1 -- estimacion optimista del valor de los candidatos restantes
public static Double heuristic(CandidatosProblem v) {
    // asume que todos los candidatos restantes aportan como mucho 5 de valoracion
    return (double)(Datos1.getNumCandidatos() - v.indice()) * 5.0;
}
\end{lstlisting}

\textbf{Nota:} una heurística perfecta ($h = \text{coste real restante}$) llevaría a A* directamente a la solución. Una heurística de cero convierte A* en BFS/Dijkstra.

% ============================================================
\section{Métodos para Backtracking (BT)}
% ============================================================

BT explora el árbol de decisiones en profundidad, podando ramas y retrocediendo cuando detecta estados inviables o peores que la mejor solución conocida.

\subsection{actions()}

Lista de \textbf{alternativas en el nodo actual}. BT las itera una a una probando cada rama. Misma semántica que en A*.

\begin{lstlisting}[style=java]
// CandidatosProblem -- contratar o no al candidato actual
public List<Boolean> actions() {
    if (indice == Datos1.getNumCandidatos()) return List.of();
    List<Boolean> lBool = List2.of(false); // siempre se puede no contratar
    Boolean c = Datos1.getSueldoMin(indice) <= this.presRest
        && this.candSel.stream()
            .noneMatch(i -> Datos1.getSonIncompatibles(this.indice(), i));
    if (c) lBool.add(true);
    return lBool;
}
\end{lstlisting}

\subsection{neighbor(action)}

Genera el \textbf{estado hijo} al tomar la acción \texttt{a}. BT desciende por este hijo. Misma semántica que en A*.

\begin{lstlisting}[style=java]
// CandidatosProblem -- avanza al siguiente candidato
public CandidatosProblem neighbor(Boolean a) {
    IntegerSet sel2 = IntegerSet.copy(candSel);
    Integer presRest2 = presRest;
    if (a) {
        sel2.add(indice);
        presRest2 = (int)(presRest2 - Datos1.getSueldoMin(indice));
    }
    return CandidatosProblem.of(indice + 1, sel2, presRest2);
}
\end{lstlisting}

\subsection{edge(action)}

Arista hacia el vecino. BT la usa para construir el \texttt{GraphPath} de la solución encontrada.

\begin{lstlisting}[style=java]
public Edge2 edge(Integer a) {
    return Edge2.of(this, this.neighbor(a), a);
}
\end{lstlisting}

\subsection{goal()}

Indica que ya no hay más decisiones. BT lo comprueba al descender para saber si ha llegado al final del árbol.

\begin{lstlisting}[style=java]
public Boolean goal() {
    return this.indice() == Datos2.getNumElementos();
}
\end{lstlisting}

\subsection{goalHasSolution()}

Si \texttt{goal()} es \texttt{true}, decide si la solución construida es válida. BT la usa para actualizar la \textbf{mejor solución conocida} (cota).

\begin{lstlisting}[style=java]
public Boolean goalHasSolution() {
    return this.indice == Datos3.N
        && this.nConsec >= 2
        && this.durAcum >= 0.0;
}
\end{lstlisting}

\subsection{isValid()}

\textbf{Poda principal de BT.} Si devuelve \texttt{false}, corta la rama entera sin explorarla. Es donde reside la mayor parte de la eficiencia de BT: cuanto antes se detecte un estado inviable, más ramas se ahorran.

\begin{lstlisting}[style=java]
public Boolean isValid() {
    return this.indice >= 0
        && this.indice <= Datos1.getNumCandidatos()
        && this.presRest <= Datos1.getPresupuestoMax();
}
\end{lstlisting}

\subsection{greedyAction() --- exclusivo BT con semilla}

Devuelve la \textbf{acción más prometedora} según un criterio voraz. BT la usa una sola vez al inicio (con \texttt{BT.ofGreedy}) para construir una primera solución completa que sirve de \textbf{cota inicial}, permitiendo podar más ramas desde el principio.

\begin{lstlisting}[style=java]
// Vertex2I -- elige el contenedor con menos hueco libre tras meter el elemento
// (estrategia "mejor ajuste": minimiza el espacio desperdiciado)
public Integer greedyAction() {
    Function<Integer, Integer> f =
        a -> capRest().get(a) - Datos2.getTamElemento(indice);
    return actions().stream()
        .min(Comparator.comparing(a -> f.apply(a)))
        .get();
}
\end{lstlisting}

% ============================================================
\section{Tabla resumen}
% ============================================================

\begin{table}[h!]
\centering
\renewcommand{\arraystretch}{1.4}
\begin{tabularx}{\textwidth}{>{\bfseries}l X c c}
\toprule
\rowcolor{tabheader}
\textbf{Método} & \textbf{Qué hace} & \textbf{A*} & \textbf{BT} \\
\midrule
\texttt{actions()} & Lista de acciones/decisiones posibles desde el estado actual & \checkmark & \checkmark \\
\texttt{neighbor(a)} & Devuelve el estado siguiente al tomar la acción \texttt{a} & \checkmark & \checkmark \\
\texttt{edge(a)} & Crea la arista hacia el vecino (wrapper de \texttt{neighbor}) & \checkmark & \checkmark \\
\texttt{goal()} & Indica si ya no quedan decisiones por tomar & \checkmark & \checkmark \\
\texttt{goalHasSolution()} & Comprueba si el estado final es una solución válida & \checkmark & \checkmark \\
\texttt{isValid()} & Poda: descarta el nodo y su subrama si devuelve \texttt{false} & \checkmark & \checkmark \\
\midrule
\textbf{Heurística} (clase externa) & Estima $h(n)$: coste optimista restante hasta la solución & \textbf{exclusivo} & --- \\
\texttt{greedyAction()} & Acción voraz para construir la cota inicial de BT & --- & \textbf{exclusivo*} \\
\bottomrule
\end{tabularx}
\caption{Métodos comunes y exclusivos de cada algoritmo. (*) Solo necesario con \texttt{BT.ofGreedy()}.}
\end{table}

\bigskip
\noindent\textbf{Conclusión:} el \texttt{Vertex} se implementa una sola vez y es compatible con ambos algoritmos. Lo que diferencia A* de BT no está dentro del vértice, sino en los elementos externos: la \textbf{heurística} para A* y \textbf{greedyAction} para BT con semilla.

\end{document}
